\section{Matrix Toolkit (Quick Reference)}

This section provides a quick summary of the key matrix definitions, operations, and formulas used throughout this booklet.

\subsection*{1. Notation and Dimensions}
An $m \times n$ matrix $A$ has $m$ rows and $n$ columns.
The entry in the $i$-th row and $j$-th column is denoted $a_{ij}$.
\[
A = \begin{pmatrix}
a_{11} & a_{12} & \cdots & a_{1n} \\
a_{21} & a_{22} & \cdots & a_{2n} \\
\vdots & \vdots & \ddots & \vdots \\
a_{m1} & a_{m2} & \cdots & a_{mn}
\end{pmatrix}
\]

\subsection*{2. Matrix Arithmetic}
\textbf{Addition/Subtraction (same dimensions only):}
\[ A \pm B \text{ adds/subtracts corresponding elements.} \]
\textbf{Scalar Multiplication:}
\[ kA \text{ multiplies every element of } A \text{ by the scalar } k. \]

\subsection*{3. Matrix Multiplication}
For the product $AB$ to be defined, the number of columns in $A$ must equal the number of rows in $B$.
\[
A_{m\times n} B_{n\times p} \to (AB)_{m\times p}
\]
Note: Matrix multiplication is \textbf{not} commutative ($AB \neq BA$ in general).

\subsection*{4. Transpose}
The transpose $A^T$ is formed by swapping rows and columns.
\[ (A^T)_{ij} = a_{ji} \]
Properties: $(A^T)^T = A$, and $(AB)^T = B^T A^T$.

\subsection*{5. Identity and Zero Matrices}
\textbf{Identity Matrix $I$:} A square matrix with 1s on the main diagonal and 0s elsewhere.
\[ IA = A \quad \text{and} \quad AI = A \]
\textbf{Zero Matrix $O$:} A matrix where all entries are 0.
\[ A + O = A \quad \text{and} \quad AO = O \]

\subsection*{6. Determinant and Inverse ($2 \times 2$)}
For a $2 \times 2$ matrix $A = \begin{pmatrix} a & b \\ c & d \end{pmatrix}$:
\begin{itemize}
    \item \textbf{Determinant:} $\det(A) = |A| = ad - bc$
    \item \textbf{Inverse:} If $\det(A) \neq 0$, the inverse exists:
    \[
    A^{-1} = \frac{1}{ad - bc} \begin{pmatrix} d & -b \\ -c & a \end{pmatrix}
    \]
\end{itemize}
Properties: $A A^{-1} = A^{-1} A = I$, and $(AB)^{-1} = B^{-1} A^{-1}$.

\subsection*{7. Solving Matrix Equations}
To solve a linear system $AX = B$:
\[
AX = B \implies X = A^{-1} B \quad \text{(if } A^{-1} \text{ exists)}
\]
(Pre-multiply both sides by $A^{-1}$).

\subsection*{8. Geometric Area via Determinants}
\textbf{Triangle Area (with one vertex at the origin $O(0,0)$):}
\[
\text{Area}(\triangle OPQ) = \frac{1}{2} |x_1 y_2 - x_2 y_1| = \frac{1}{2} \left| \det \begin{pmatrix} x_1 & y_1 \\ x_2 & y_2 \end{pmatrix} \right|
\]
\textbf{General Triangle Area (vertices $P_1, P_2, P_3$):}
\[
\text{Area}(\triangle P_1 P_2 P_3) = \frac{1}{2} \left| \det \begin{pmatrix} x_1 & y_1 & 1 \\ x_2 & y_2 & 1 \\ x_3 & y_3 & 1 \end{pmatrix} \right| = \frac{1}{2} \left| \det \begin{pmatrix} x_2 - x_1 & y_2 - y_1 \\ x_3 - x_1 & y_3 - y_1 \end{pmatrix} \right|
\]
\textbf{Collinearity:} Points $P_1, P_2, P_3$ are collinear if and only if $\det = 0$.
